% GENERATED FILE. Edit canonical lessons, then run npm run generate:books.
% canonical-source-hash: fnv1a64:19ae41547c10a3ac
% canonical-lessons: GU-C442-samudra, GU-C442-laher, GU-C442-nav, GU-C442-jal, GU-C442-dolvu

\chapter{Waves and Nets}
\label{ch:gu-waves-and-nets}
\hlchaptermodality{442}

\begin{chapteropening}
\textbf{By the end of this chapter:} \emph{I can talk about the sea and boats.}

Complete the last lesson of chapter 442: I can talk about the sea and boats.
\end{chapteropening}

\section[\texorpdfstring{\gu{સમુદ્ર} (samudra)}{સમુદ્ર (samudra)}]{\gu{સમુદ્ર} (samudra) — the ocean; the sea (formal)}

\label{lesson:GU-C442-samudra}



\begin{quote}
Before the new one: say the Gujarati for a riverbank, then the Gujarati for to sink in.
\end{quote}

\subsection*{You'll want to know: \gu{સમુદ્ર}}
\textbf{\gu{સમુદ્ર}} — \emph{samudra} — \textquotedblleft{}the ocean; the sea (formal)\textquotedblright{}.

\begin{quote}
\gu{સમુદ્રનું} \gu{પાણી} \gu{ખારું} \gu{છે}.
\end{quote}

\begin{quote}
\emph{samudranũ pāṇī khārũ che.}
\end{quote}

\begin{quote}
The water of the ocean is salty.
\end{quote}

A masculine noun.

\begin{cousinweb}
\textbf{\gu{સમુદ્ર}} (\emph{samudra}) is a tatsama, Sanskrit \emph{samudra}, \textquotedblleft{}a gathering together of waters\textquotedblright{}, the sea (Monier-Williams); speech says \gu{દરિયો}.
\end{cousinweb}

\begin{grammarlens}[title={What you've built}]
One more word for getting around.
\end{grammarlens}

\subsection*{Your turn}
\begin{itemize}
\raggedright
  \item \emph{Say it:} \emph{samudra}
  \item \emph{Say it:} \emph{samudra}, once more
  \item \emph{Say it:} \emph{samudranũ pāṇī khārũ che.} — The water of the ocean is salty.
  \item \emph{Recall:} say the Gujarati for a riverbank, then the Gujarati for to sink in, then say \emph{samudra} again
\end{itemize}

\begin{tcolorbox}[breakable,colback=teal!4,colframe=teal!35!black,title={Before you move on}]
What does \gu{સમુદ્ર} mean? (\textquotedblleft{}The ocean; the sea (formal)\textquotedblright{}.) Say it once more.
\end{tcolorbox}
\section[\texorpdfstring{\gu{લહેર} (laher)}{લહેર (laher)}]{\gu{લહેર} (laher) — a wave; a ripple}

\label{lesson:GU-C442-laher}



\begin{quote}
Before the new one: say the Gujarati for to sink in, then the Gujarati for the ocean.
\end{quote}

\subsection*{You'll want to know: \gu{લહેર}}
\textbf{\gu{લહેર}} — \emph{laher} — \textquotedblleft{}a wave; a ripple\textquotedblright{}.

\begin{quote}
\gu{દરિયાની} \gu{લહેરો} \gu{કિનારા} \gu{પર} \gu{આવે} \gu{છે}.
\end{quote}

\begin{quote}
\emph{dariyānī lahero kinārā par āve che.}
\end{quote}

\begin{quote}
The sea's waves come up on the shore.
\end{quote}

A feminine noun.

\begin{cousinweb}
\textbf{\gu{લહેર}} (\emph{laher}) is a tadbhava of Sanskrit \emph{laharī}, \textquotedblleft{}a wave\textquotedblright{}; CDIAL lists the Gujarati word (CDIAL 10999).
\end{cousinweb}

\begin{grammarlens}[title={What you've built}]
One more word for getting around.
\end{grammarlens}

\subsection*{Your turn}
\begin{itemize}
\raggedright
  \item \emph{Say it:} \emph{laher}
  \item \emph{Say it:} \emph{laher}, once more
  \item \emph{Say it:} \emph{dariyānī lahero kinārā par āve che.} — The sea's waves come up on the shore.
  \item \emph{Recall:} say the Gujarati for to sink in, then the Gujarati for the ocean, then say \emph{laher} again
\end{itemize}

\begin{tcolorbox}[breakable,colback=teal!4,colframe=teal!35!black,title={Before you move on}]
What does \gu{લહેર} mean? (\textquotedblleft{}A wave; a ripple\textquotedblright{}.) Say it once more.
\end{tcolorbox}
\section[\texorpdfstring{\gu{નાવ} (nāv)}{નાવ (nāv)}]{\gu{નાવ} (nāv) — a river boat}

\label{lesson:GU-C442-nav}



\begin{quote}
Before the new one: say the Gujarati for the ocean, then the Gujarati for a wave.
\end{quote}

\subsection*{You'll want to know: \gu{નાવ}}
\textbf{\gu{નાવ}} — \emph{nāv} — \textquotedblleft{}a river boat\textquotedblright{}.

\begin{quote}
\gu{નાવમાં} \gu{બેસીને} \gu{અમે} \gu{નદી} \gu{પાર} \gu{કરી}.
\end{quote}

\begin{quote}
\emph{nāvmā̃ besīne ame nadī pār karī.}
\end{quote}

\begin{quote}
We crossed the river in a boat.
\end{quote}

A feminine noun.

\begin{cousinweb}
\textbf{\gu{નાવ}} (\emph{nāv}) is a tadbhava of Sanskrit \emph{nāvā}, \textquotedblleft{}a boat\textquotedblright{}; CDIAL lists the Gujarati word (CDIAL 7081). Songs and stories prefer it to \gu{હોડી}.
\end{cousinweb}

\begin{grammarlens}[title={What you've built}]
One more word for getting around.
\end{grammarlens}

\subsection*{Your turn}
\begin{itemize}
\raggedright
  \item \emph{Say it:} \emph{nāv}
  \item \emph{Say it:} \emph{nāv}, once more
  \item \emph{Say it:} \emph{nāvmā̃ besīne ame nadī pār karī.} — We crossed the river in a boat.
  \item \emph{Recall:} say the Gujarati for the ocean, then the Gujarati for a wave, then say \emph{nāv} again
\end{itemize}

\begin{tcolorbox}[breakable,colback=teal!4,colframe=teal!35!black,title={Before you move on}]
What does \gu{નાવ} mean? (\textquotedblleft{}A river boat\textquotedblright{}.) Say it once more.
\end{tcolorbox}
\section[\texorpdfstring{\gu{જાળ} (jāḷ)}{જાળ (jāḷ)}]{\gu{જાળ} (jāḷ) — a fishing net}

\label{lesson:GU-C442-jal}



\begin{quote}
Before the new one: say the Gujarati for a wave, then the Gujarati for a river boat.
\end{quote}

\subsection*{You'll want to know: \gu{જાળ}}
\textbf{\gu{જાળ}} — \emph{jāḷ} — \textquotedblleft{}a fishing net\textquotedblright{}.

\begin{quote}
\gu{સવારે} \gu{હોડીમાંથી} \gu{તેમણે} \gu{દરિયામાં} \gu{જાળ} \gu{નાખી}.
\end{quote}

\begin{quote}
\emph{savāre hoḍīmā̃thī temṇe dariyāmā̃ jāḷ nākhī.}
\end{quote}

\begin{quote}
In the morning they cast a net into the sea from the boat.
\end{quote}

A feminine noun.

\begin{cousinweb}
\textbf{\gu{જાળ}} (\emph{jāḷ}) is a tadbhava of Sanskrit \emph{jāla}, \textquotedblleft{}a net\textquotedblright{}; CDIAL lists the Gujarati word (CDIAL 5213).
\end{cousinweb}

\begin{grammarlens}[title={What you've built}]
One more word for getting around.
\end{grammarlens}

\subsection*{Your turn}
\begin{itemize}
\raggedright
  \item \emph{Say it:} \emph{jāḷ}
  \item \emph{Say it:} \emph{jāḷ}, once more
  \item \emph{Say it:} \emph{savāre hoḍīmā̃thī temṇe dariyāmā̃ jāḷ nākhī.} — In the morning they cast a net into the sea from the boat.
  \item \emph{Recall:} say the Gujarati for a wave, then the Gujarati for a river boat, then say \emph{jāḷ} again
\end{itemize}

\begin{tcolorbox}[breakable,colback=teal!4,colframe=teal!35!black,title={Before you move on}]
What does \gu{જાળ} mean? (\textquotedblleft{}A fishing net\textquotedblright{}.) Say it once more.
\end{tcolorbox}
\section[\texorpdfstring{\gu{ડોલવું} (ḍolvũ)}{ડોલવું (ḍolvũ)}]{\gu{ડોલવું} (ḍolvũ) — to rock gently; to sway (on water)}

\label{lesson:GU-C442-dolvu}



\begin{quote}
Before the new one: say the Gujarati for a river boat, then the Gujarati for a fishing net.
\end{quote}

\subsection*{You'll want to know: \gu{ડોલવું}}
\textbf{\gu{ડોલવું}} — \emph{ḍolvũ} — \textquotedblleft{}to rock gently; to sway (on water)\textquotedblright{}.

\begin{quote}
\gu{લહેરોમાં} \gu{હોડી} \gu{ડોલે} \gu{છે}.
\end{quote}

\begin{quote}
\emph{laheromā̃ hoḍī ḍole che.}
\end{quote}

\begin{quote}
The boat rocks on the waves.
\end{quote}

Past (it …): \gu{તે} \gu{હોડી} \gu{ડોલી}, \emph{te hoḍī ḍolī}.

\begin{cousinweb}
\textbf{\gu{ડોલવું}} (\emph{ḍolvũ}) is a tadbhava of Sanskrit \emph{dōlāyatē}, \textquotedblleft{}swings to and fro\textquotedblright{}; CDIAL lists the Gujarati verb, \textquotedblleft{}to dangle, to oscillate\textquotedblright{} (CDIAL 6585).
\end{cousinweb}

\begin{grammarlens}[title={What you've built}]
One more word for getting around.
\end{grammarlens}

\subsection*{Your turn}
\begin{itemize}
\raggedright
  \item \emph{Say it:} \emph{ḍolvũ}
  \item \emph{Say it:} \emph{ḍolvũ}, once more
  \item \emph{Say it:} \emph{laheromā̃ hoḍī ḍole che.} — The boat rocks on the waves.
  \item \emph{Recall:} say the Gujarati for a river boat, then the Gujarati for a fishing net, then say \emph{ḍolvũ} again
\end{itemize}

\begin{tcolorbox}[breakable,colback=teal!4,colframe=teal!35!black,title={Before you move on}]
What does \gu{ડોલવું} mean? (\textquotedblleft{}To rock gently; to sway (on water)\textquotedblright{}.) Say it once more.
\end{tcolorbox}
